
The methodology operationalizes signature detection through three components: multi-dimensional performance measurement, statistical clustering, and effect size quantification.

\subsubsection{Overview}

Given N models aligned with different methods, the pipeline:

\begin{enumerate}
\item Measures each model across three quality dimensions on standardized tasks (HumanEval+)
\item Projects performance vectors into 3D space via PCA
\item Clusters models by alignment method and computes intercluster distance
\item Validates signature existence via Cohen's d effect size (threshold: d > 1.$5\sigma )$
\end{enumerate}

\subsubsection{Multi-Dimensional Performance Measurement}

A performance signature is defined as a vector \textbf{s}$_{m}$ = (c\_m, x\_m, e\_m) for model m, where:

\begin{itemize}
\item \textbf{Correctness} c\_m: Pass@k rate on HumanEval+ test suites (execution success)
\item \textbf{Complexity} x\_m: Cyclomatic complexity averaged across generated code
\item \textbf{Efficiency} e\_m: Runtime (milliseconds) and memory usage (MB) averaged across successful executions
\end{itemize}

These dimensions are independent (correctness does not determine complexity or efficiency), objective (no human annotation required), scalable (automated evaluation), and relevant to practitioners.

\#\#\#\# Dimension Rationale

\textbf{Correctness} is the standard evaluation metric. Pass@k measures functional correctness via test execution. HumanEval+ provides extended test suites (3-$5\times$ more tests than HumanEval) to reduce variance. Execution-based alignment methods explicitly optimize this dimension via pass/fail feedback.

\textbf{Complexity} measures code structure independent of correctness. A model can achieve 100\% pass@k with simple or convoluted implementations. Cyclomatic complexity counts independent execution paths through code. Complexity is not measured during alignment training for execution-based methods, representing an implicit dimension where signatures may emerge.

\textbf{Efficiency} captures computational cost. Code can pass tests quickly or slowly, use memory sparingly or excessively. Runtime is measured via cProfile, peak memory via memory\_profiler, averaging across successful executions. Like complexity, efficiency is typically unmeasured during alignment.

\#\#\#\# Measurement Protocol

For each model m and task t in HumanEval+:

\begin{enumerate}
\item Generate k=10 code samples at temperature T=0.8
\item Execute each sample against HumanEval+ test suite, recording pass/fail
\item Compute correctness $c_{m,t}$ = (passed samples)/k
\item Analyze each sample's AST for cyclomatic complexity, average across samples
\item Profile successful executions for runtime and memory, average across passed samples
\end{enumerate}

Aggregate across tasks to obtain model-level signatures.

\subsubsection{Statistical Clustering and Signature Detection}

\#\#\#\# PCA-Based Dimensionality Reduction

Raw performance vectors lie in 3D space defined by correctness, complexity, and efficiency axes. PCA identifies dominant variance directions. If alignment methods create signatures, models should separate along principal components by alignment method type.

\#\#\#\# k-Means Clustering

k-means clustering with k=3 clusters (execution-based, preference-based, baseline) tests whether models cluster by alignment method. Alignment purity measures clustering quality:

Purity = (1/N) $\Sigma$ max\_method |{m $\in C_j$ : m uses method}|

If models cluster by alignment method (not architecture), purity approaches 1.0.

\#\#\#\# Cohen's d Effect Size

Cohen's d quantifies intercluster separation:

$d = \bar{D}_{\text{inter}} / \sqrt{\sigma^{2}_{\text{intra}}}$

where $\bar{D}_{\text{inter}}$ is mean pairwise distance between cluster centroids and $\sigma^{2}_{\text{intra}}$ is pooled within-cluster variance. The pre-registered threshold is $d > 1.5$ (very large effect size), ensuring detected signatures are robust rather than marginal statistical artifacts.

\subsubsection{Mechanistic Validation}

Detecting signatures (Cohen's d > 1.5) establishes that alignment methods create differences. To validate why, mechanistic predictions test dimension-specific patterns:

\textbf{M1 (Execution Dominance)}: If execution feedback optimizes correctness, then execution-based models should rank in the top 15\% on correctness across all models.

\textbf{M2 (Preference Balance)}: If preference feedback captures holistic quality, then preference-based models should achieve balanced top-30\% performance across all dimensions.

These mechanistic checks validate the causal chain: feedback signal type $\rightarrow$ optimization priority $\rightarrow$ signature pattern.

\subsubsection{Design Justification}

Three dimensions capture orthogonal quality aspects without requiring subjective annotation. Cohen's d > 1.$5\sigma$ provides a conservative threshold (80\%+ non-overlap between distributions). HumanEval+ offers standardization (164 function-level tasks, extended test suites, widely used for alignment evaluation).

\subsubsection{Limitations}

Complexity and efficiency metrics are proxies for code quality. Testing 3-4 models (vs. planned 6-8) reduces statistical power. The proof-of-concept uses simulated data for methodology validation, with real-model validation required for generalization. Model scale confounds (comparing 2.7B vs 350M models) may conflate alignment effects with capacity effects.
